\documentclass[12pt,a4paper]{article}
\usepackage{amsmath}
\usepackage{amssymb}  % Provides \mathbb and other symbols
\usepackage{graphicx}
\usepackage{subcaption}
\usepackage{enumitem}
\usepackage{listings}
\usepackage{hyperref}
\usepackage{ulem}
\usepackage{xcolor}
\usepackage{amsthm}
\usepackage{tcolorbox}  % For colored boxes
\usepackage[table]{xcolor} % For coloring table cells
\usepackage{array}         % For better table alignment
\usepackage{mdframed}
\usepackage{cancel}
\usepackage[margin=2.5cm]{geometry}
\setlength{\parindent}{0pt}
\setlength{\parskip}{0.8em}

\begin{document}
\begin{titlepage}
\centering
\vspace*{\fill}
{\Huge Analysis III} \\[0.5cm]
{\Huge Home Work 3}\\[1.5cm]
{\Large Student: 26-712-224}\\[0.5cm]
{\Large \today}
\vspace*{\fill}
\end{titlepage}

\newpage
\textbf{Exercise 1 (Properties of $L^1$, 1+1+1+1+1 P.)}
For $f:\mathbb{R}^n\to\mathbb{R}$, define the positive part $f_+:=\max\{0,f\}$ and negative part $f_-:=\max\{0,-f\}$ of $f$ respectively. The goal of this exercise is to show that
\begin{align*}
f\in L^1(\mathbb{R}^n;\mathbb{R}) \iff f_{\pm}\in L^1(\mathbb{R}^n;\mathbb{R}).
\end{align*}
To do this, follow the procedure outlined below.
\begin{enumerate}
\item[(a)] Prove the ``$\Leftarrow$'' direction.
\item[(b)] To prove ``$\Rightarrow$'', let $f\in L^1(\mathbb{R}^n;\mathbb{R})$. By definition, there are $g,h\in\mathcal{L}^{\mathrm{inc}}$ with $f=g-h$. Write down the properties of the sequences $(g_k)_{k\in\mathbb{N}}$ and $(h_k)_{k\in\mathbb{N}}$ obtained by applying the definition of $\mathcal{L}^{\mathrm{inc}}$ on $g$ and $h$ respectively.
\item[(c)] Justify briefly that $(r_k)_{k\in\mathbb{N}}$, defined by $r_k=2\min\{g_k,h_k\}$, is an increasing sequence of step functions (you do not need to give a full proof). Prove that the sequence
\begin{align*}
\left(\int r_k(x)\,dx\right)_{k\in\mathbb{N}}
\end{align*}
is bounded. \textit{Hint:} You may use without proof that $2\min\{a,b\}\leq a+b$.
\item[(d)] Prove that
\begin{align*}
r_k=g_k+h_k-|g_k-h_k|
\end{align*}
and find the limit as $k\to+\infty$. Using that $\mathcal{L}^{\mathrm{inc}}\subset L^1(\mathbb{R}^n;\mathbb{R})$ and that $L^1(\mathbb{R}^n;\mathbb{R})$ is a vector space, deduce that $|f|\in L^1(\mathbb{R}^n;\mathbb{R})$.
\item[(e)] Express $f_+$ and $f_-$ in terms of $f$ and $|f|$ and hence deduce that $f_{\pm}\in L^1(\mathbb{R}^n;\mathbb{R})$.
\end{enumerate}


\vspace{1cm}
\noindent
Answer:\\
\\
(a) From the definition of $f_+$ and $f_-$ it follows that: $f = f_+ - f_-$.  But we are given that both: $f_-, f_+ \in L^1$ and as $L^1$
is a vector space we must also have that $f \in L^1$.\\
\\
(b) From what is already given, we can conclude that there exist sequences $(g_k), (h_k) \in T^{inc}$ such that:
\begin{itemize}
\item $h_k \to h$ and $g_k \to g$ - both pointwise a.e.
\item $(\int h_k)$ and $(\int g_k)$both converge.
\end{itemize}
These are precisely the properties given by the definition of $L^{inc}$.\\
\\
(c) The minimum of two step functions is a step function.  By the hint and monotonicity,
\begin{align*}
\int r_1 \leq \int r_k
\leq \int g_k+\int h_k.
\end{align*}
As $g_k$ and $h_k$ are both increasing means $r_k$ also increasing.  The right-hand side 
converges, so $(\int r_k)$ is bounded. Since it is also increasing, it must converge.\\
\\
(d) The identity $2\min\{a,b\} = a+b-|a-b|$ gives:
\begin{align*}
r_k=g_k+h_k-|g_k-h_k|.
\end{align*}
Outside the union of the two null exceptional sets,
\begin{align*}
r_k\longrightarrow r:=g+h-|g-h|=g+h-|f|.
\end{align*}
Part (c) therefore gives $r \in L^{inc} \subset L^1$. Since $g,h \in L^1$ and $L^1$ is a vector space,
\begin{align*}
|f|=g+h-r\in L^1.
\end{align*}
\\
(e) We can write: $f_- = \frac{|f| - f}{2}$ and $f_+ = \frac{|f| + f}{2}$.  As $f \in L^1 \implies |f| \in L^1$.  Also, we know that $L^1$ is
a vector space.  This means that, using the above two equations, that $f_+$ and $f_-$ are both $L^1$. 
\newpage
\textbf{Exercise 2 (Integrability, 1+1+1+1+1 P.)}
For the following functions, decide whether they are Lebesgue integrable, Riemann integrable, both or neither. You do not need to compute the integrals. Cite Theorems you are using from the script. You may also use exercise 1 above.
\begin{enumerate}
\item[(a)] $f_1(x)=\pi$ on $\mathbb{R}$
\item[(b)] $f_2(x)=-e^{-x}$ on $\mathbb{R}$
\item[(c)] $f_3(x)=e^{-|x|}$ on $\mathbb{R}$
\item[(d)] $f_4(x)=\chi_{\mathbb{R}\setminus\mathbb{Q}}(x)-\chi_{\mathbb{Q}}(x)$ on $[0;1]$
\item[(e)] $f_5(x)=x\sin(1/x)$, $f_5(0)=0$ on $[-1;1]$
\end{enumerate}

\vspace{1cm}
\noindent
Answer:\\
\\
(a) This function is not Lebesgue Integrable.  If $f_1$ were Lebesgue Integrable, monotonicity (Theorem 11.2.14(2)) would give:
\begin{align*}
\int_{\mathbb{R}} f_1 \geq \int_{[-N,N]} \pi = 2\pi N
\qquad \text{for every } N.
\end{align*}
contradicting its finite integral.\\
\\
The Riemann Integral also doesn't exists as the integral over $[-N,N]$ tends to infinity as $N \to \infty$.\\
\\
(b) This function is not Lebesgue Integrable. If $f_2$ were Lebesgue Integrable, $|f_2|$ would also be Lebesque Integrable (Theorem 11.2.14(3)). Monotonicity would then give:
\begin{align*}
\int_{\mathbb{R}} |f_2| \geq \int_{[-N,0]} e^{-x} = e^N - 1.
\end{align*}
which does not converge.  This contradictin proves the point.\\
\\
Not improperly Riemann Integral: its integral over $[-N,0]$ is $1-e^N$, tending to $-\infty$.\\
\\
(c) This function is Lebesgue Integrable.  Its improper Riemann integral converges absolutely, so Theorem 11.2.17(2) applies.\\
\\
The Reimann Integral also converges (improperly) to 2.\\
\\
(d) This function is Lebesgue Integrable.  The rationals form a null set and using theorem 11.2.15 the intregal will
be the same as integrating $f=1$ on $[0,1]$.\\
\\
Not Riemann Integrable - every subinterval contains rationals and irrationals, giving lower sums $-1$ and upper sums $1$.\\
\\
(e) The function is Lebesgue Integrable.  It is continuous, including at 0 since $|x \sin(1/x)| \le |x|$. Then apply 
Theorem 11.2.16(2).

\newpage
\noindent
\textbf{Exercise 5 (Not Lebesgue-integrable function)}
Show that the integral
\begin{align*}
\int_1^{+\infty} \frac{\sin(x)}{x}\,dx
\end{align*}
exists as a Riemann integral, but not as a Lebesgue integral. Why is this not in contradiction with Theorem 11.2.17(2)? \textit{Hint:} Integration by parts.

\vspace{1cm}
\noindent
Answer:\\
\\
For such Reimann integrals one usually takes an upper limit $R$ - which we later lets tend to infinity.\\
\\
Integration by parts gives, for $R > 1$:
\begin{align*}
\int_1^R\frac{\sin(x)}{x} \; dx
=\cos 1-\frac{\cos(R)}{R}-\int_1^R\frac{\cos(x)}{x^2} \; dx.
\end{align*}
The boundary term tends to zero; the last integral converges absolutely since $|\cos(x)|/x^2 \le 1/x^2$. Hence the improper Riemann Integral exists.\\
\\
\textcolor{red}{Sorry - not enough time to complete this.}
\end{document}